\section*{ЗАКЛЮЧЕНИЕ}
\addcontentsline{toc}{section}{ЗАКЛЮЧЕНИЕ}

В ходе прохождения проектной (учебной) практики был успешно реализован комплексный медиа-проект «ЭкоМедиа», направленный на просвещение молодежи в области экологии человека и цифровой гигиены. Выполнение работы потребовало интеграции теоретических знаний об актуальных экологических угрозах (микропластик, водопроводная вода, влияние гаджетов на физиологию сна) с современными информационными технологиями.

Поставленная в начале проекта цель — создание инновационного медиа-ресурса и интерактивного инструмента для формирования здоровых привычек в условиях городской среды — была полностью достигнута.

В рамках \textbf{базовой части задания} был спроектирован и сверстан адаптивный веб-сайт. Использование современного визуального стиля Glassmorphism и темной темы оформления (Dark Mode) не только сделало ресурс визуально привлекательным для целевой аудитории, но и на практике реализовало один из принципов цифровой гигиены — снижение нагрузки на зрительный аппарат пользователя. Команда успешно освоила инструменты совместной разработки, настроив инфраструктуру на базе Git и GitHub, что обеспечило прозрачность и контроль версионирования исходного кода.

В рамках \textbf{вариативной части задания} (Вариант 2) был разработан полнофункциональный Telegram-бот. Применение языка программирования Python и библиотеки \texttt{pyTelegramBotAPI} позволило создать стабильный программный продукт, реализующий сложную логику взаимодействия с пользователем на основе конечного автомата (FSM). Интерактивные модули эко-викторины и тестирования уровня цифровой гигиены показали высокую эффективность вовлечения пользователей во время альфа-тестирования на фокус-группе.

Таким образом, все поставленные задачи были решены в полном объеме. Практическая значимость работы заключается в том, что разработанный программно-информационный комплекс представляет собой готовое к эксплуатации решение (MVP), которое может быть развернуто на сервере и использовано образовательными или некоммерческими экологическими организациями для проведения массовых просветительских кампаний.

В процессе выполнения проекта авторы получили ценный опыт в области веб-разработки (HTML5, CSS3), программирования (Python), проектирования пользовательских интерфейсов (UI/UX) и работы с системами контроля версий, что заложило прочный фундамент для дальнейшего профессионального роста в сфере информационных технологий.
